\subsection{The transverse equation}\label{cc:uf:sec:transverse}
Each oscillatory correction solves a forced seed equation forward from
zero initial data at each label and phase. Its forcing vanishes before
activation, so the correction does too. From activation onward we measure
its size relative to the amplitude of the seed at the origin, using the
weight $\gamma$.

We call
\begin{equation}\label{cc:uf:hier:eq:trans}
 V_t+MV+\kappa\partial_s\mathsf p=f,\qquad
 \kappa\cdot V=0,\qquad \av V=\av f=0,
\end{equation}
the \emph{transverse equation}, because its unknown $V$ is perpendicular
to the wave vector $\kappa$. The pressure term keeps $V$
perpendicular to $\kappa$: differentiating $\kappa\cdot V=0$
and using $\kappa_t=-M^T\kappa$ in \eqref{cc:uf:hier:eq:trans} gives
\begin{equation}\label{cc:uf:hier:eq:pressure-from-trans}
 |\kappa|^2\partial_s\mathsf p=\kappa\cdot f-2\kappa\cdot MV.
\end{equation}
The transverse equation is therefore the linear ordinary differential
equation
\begin{equation}\label{cc:uf:hier:eq:trans-ode}
 V_t=-(\Id-2P_\kappa)MV+(\Id-P_\kappa)f
\end{equation}
at each fixed $(a,s)$; for $f=0$ it is the seed equation
\eqref{cc:uf:step:eq:seed}. For
every source we take its solution with $V(0,a,s)=0$, and $\mathsf p$ is
the zero-mean primitive in $s$ of
\eqref{cc:uf:hier:eq:pressure-from-trans}. For smooth $f$ the solution
is unique and smooth, and $V$ satisfies the two constraints in
\eqref{cc:uf:hier:eq:trans}. Both $V$ and $\mathsf p$ vanish at every
label where $f$ vanishes for all $(t,s)$. If $f$ vanishes for $t<t_0$,
then $V=0$ for $t\le t_0$.

The plane $\kappa^\perp$ varies with the label. To compare the equation
on the envelope with the one at the origin and to differentiate it in
the labels, we express $V$ in coordinates in a fixed $\R^2$.
Choose a constant $3\times2$ matrix with orthonormal columns spanning
$k_*^\perp$, let $A_\perp$ be $D_aX$ times this matrix, and let
$A_\perp^\dagger=(A_\perp^TA_\perp)^{-1}A_\perp^T$.
The columns of $A_\perp$ span $\kappa^\perp$, since
$\kappa=(D_aX)^{-T}k_*$, and $A_\perp^\dagger$ takes a vector to its
coordinates in those columns after orthogonal projection onto their
span: $A_\perp^\dagger A_\perp=\Id_{\R^2}$ and
$A_\perp A_\perp^\dagger=\Id-P_\kappa$. Write $V=A_\perp v$ for a
transverse velocity. Since $\partial_tA_\perp=MA_\perp$, the equation
\eqref{cc:uf:hier:eq:trans-ode} becomes
\begin{equation}\label{cc:uf:tower:eq:coordinate}
 v_t+\mathcal M_\perp v=\psi,\qquad
 \mathcal M_\perp=2A_\perp^\dagger MA_\perp,\qquad
 \psi=A_\perp^\dagger f.
\end{equation}

\begin{lemma}\label{cc:uf:tower:lem:inverse}
For $t_0\le t\le T$ and $|a|\le2\ell_*$, the seed satisfies
\begin{equation}\label{cc:uf:tower:eq:near-origin}
 |\mathbf a(t,a)|\le2|\mathbf a(t,0)|,\qquad
 (\kappa\cdot M\mathbf a)(t,a)>0 .
\end{equation}
There is an absolute exponent $C$ such that the following holds at every
order $n\ge0$.
\begin{enumerate}[label=(\roman*),leftmargin=*]
\item Let $V$ and $f$ be smooth fields on $[t_0,T]$ that satisfy
\eqref{cc:uf:hier:eq:trans-ode}, with $V$ perpendicular to $\kappa$, and
suppose that $V(t_0)$ and $f$ vanish at the labels outside $\supp g$.
Then
\begin{equation}\label{cc:uf:tower:eq:seed-norm}
 \|V\|_{[n],\gamma}\le C_n\mathsf D^C
 \bigl(\gamma(t_0)^{-1}\|V(t_0)\|_{[n]}+\|f\|_{[n],\gamma}\bigr).
\end{equation}
\item Let $f$ be smooth on $[0,T]$ with zero phase mean, and suppose
that $f$ vanishes for $t<t_0$ and at the labels outside $\supp g$.
The solution of \eqref{cc:uf:hier:eq:trans} with $V(0)=0$ and its
pressure satisfy
\begin{equation}\label{cc:uf:hier:eq:trans-bound}
 \|V\|_{[n],\gamma}+\lambda\|\mathsf p\|_{[n],\gamma}
 \le C_n\mathsf D^C\|f\|_{[n],\gamma}.
\end{equation}
\item On $[t_0,T]$ the leading oscillation and its pressure satisfy
\begin{equation}\label{cc:uf:hier:eq:size}
 \|A_1\|_{[n],\gamma}+\lambda\|\mathsf p[A_1]\|_{[n],\gamma}
 \le C_n\mathsf D^C\lambda^{-5/4}.
\end{equation}
\end{enumerate}
\end{lemma}
\begin{proof}
By \eqref{cc:uf:tower:eq:parent-bound},
$\mathsf D^{-1}|w|\le|A_\perp w|\le\mathsf D|w|$ for $w\in\R^2$. Hence
$|A_\perp|\le\mathsf D$ and $|A_\perp^\dagger|\le\mathsf D$, and the
matrix $K_\perp=A_\perp^TA_\perp$ satisfies
$\sup|K_\perp^{-1}|\le\mathsf D^2$. By
Proposition~\ref{cc:uf:trunc:prop:parent} and
Lemma~\ref{cc:uf:trunc:lem:coefficients}(a),(b), the coefficients
$A_\perp$, $A_\perp^\dagger=K_\perp^{-1}A_\perp^T$ and $\mathcal M_\perp$
satisfy \eqref{cc:uf:trunc:eq:power}.

In the coordinates $V=A_\perp v$, the seed equation at the label $a$ is
$v_t+\mathcal M_\perp(t,a)v=0$. With the bounds for $A_\perp$ and
$A_\perp^\dagger$, the first condition
in \eqref{cc:uf:tower:eq:wave-origin} gives
$|v(t)|/\gamma(t)\le\mathsf D^3|v(t')|/\gamma(t')$ for its solutions at
$a=0$. By the bound for the first label
derivatives of $\mathcal M_\perp$ and the choice of $\ell_*$,
\[
 \sup_{t_0\le t\le T}
 |\mathcal M_\perp(t,a)-\mathcal M_\perp(t,0)|
 \le C\mathsf D^C|a|\le\mathsf D^{-3}\qquad(|a|\le2\ell_*).
\]
Let $v$ solve the equation at such a label $a$ on $[t',T]$. It solves
the equation at the origin with
the source $-(\mathcal M_\perp(\cdot,a)-\mathcal M_\perp(\cdot,0))v$, so
variation of constants gives
\[
 \frac{|v(t)|}{\gamma(t)}
 \le\mathsf D^3\frac{|v(t')|}{\gamma(t')}
 +\int_{t'}^t\frac{|v(t'')|}{\gamma(t'')}\dd t''.
\]
The weights cancel because the bound at the origin is a ratio of values
of $\gamma$. Since $t-t'\le T\le1$, Gronwall's inequality gives
$|v(t)|/\gamma(t)\le e\mathsf D^3|v(t')|/\gamma(t')$. For a solution of
\eqref{cc:uf:tower:eq:coordinate} on $[t_0,T]$, variation of constants
therefore gives, at each phase and each label with $|a|\le2\ell_*$,
\begin{equation}\label{cc:uf:tower:eq:weighted-forward}
 \frac{|v(t,a)|}{\gamma(t)}
 \le e\mathsf D^3\left(
 \frac{|v(t_0,a)|}{\gamma(t_0)}
 +\int_{t_0}^t\frac{|\psi(t',a)|}{\gamma(t')}\,dt'\right).
\end{equation}

To compare the seed at $a$ with its value at the origin, write
$\mathbf a=A_\perp v_a$, with $v_a=A_\perp^\dagger\mathbf a$ at $a$. The
difference $v_a-v_0$ solves the equation at the origin with the source
$-(\mathcal M_\perp(\cdot,a)-\mathcal M_\perp(\cdot,0))v_a$. Its value at
$t_0$ is $(A_\perp^\dagger(t_0,a)-A_\perp^\dagger(t_0,0))\mathbf a(t_0,0)$,
by \eqref{cc:uf:tower:eq:start} and $A_\perp^\dagger P_\kappa=0$, and
$|\mathbf a(t_0,0)|=\gamma(t_0)\lambda^{-3/4}$. With the bound just
proved for $v_a$, variation of constants gives
$|v_a(t)-v_0(t)|\le C\mathsf D^C|a|\gamma(t)\lambda^{-3/4}$. Since
$\gamma(t)\lambda^{-3/4}=|\mathbf a(t,0)|$, it follows that
\[
 |\mathbf a(t,a)-\mathbf a(t,0)|
 \le C\mathsf D^C|a|\,|\mathbf a(t,0)|,
 \qquad t_0\le t\le T.
\]
The label derivatives also give
$|\kappa(t,a)-\kappa(t,0)|\le C\mathsf D^C|a||\kappa(t,0)|$
and $|M(t,a)-M(t,0)|\le C\mathsf D^C|a|$. Since $|M|\le\mathsf D$, the
three bounds give
\[
 \bigl|(\kappa\cdot M\mathbf a)(t,a)-(\kappa\cdot M\mathbf a)(t,0)\bigr|
 \le C\mathsf D^C|a|\,|\kappa(t,0)|\,|\mathbf a(t,0)| ,
\]
while $(\kappa\cdot M\mathbf a)(t,0)
=\nu(t,0)|\kappa(t,0)||\mathbf a(t,0)|
\ge\mathsf D^{-1}|\kappa(t,0)||\mathbf a(t,0)|$, by the second condition
in \eqref{cc:uf:tower:eq:wave-origin}. None of these exponents depends
on the exponent of $\ell_*$, so the choice of $\ell_*$ proves
\eqref{cc:uf:tower:eq:near-origin}.

Let $V$ and $f$ be as in (i). The solution vanishes at the labels
outside $\supp g$, where its initial value and source vanish, and
$\supp g$ lies in $\{|a|\le2\ell_*\}$. Taking the supremum in
\eqref{cc:uf:tower:eq:weighted-forward} gives, because
$t-t_0\le T\le1$,
\[
 \|v\|_{[0],\gamma}\le e\mathsf D^3\bigl(\gamma(t_0)^{-1}\|v(t_0)\|_{[0]}
       +\|\psi\|_{[0],\gamma}\bigr).
\]
Since $\supp g$ is closed, the label and phase derivatives of $v$ and
$\psi$ also vanish outside it. So this is the zero-order
estimate of Lemma~\ref{cc:uf:trunc:lem:lifting} on $[t_0,T]$, for the
initial value and the source together, with the one coefficient
$\mathcal M_\perp$, which enters through the identity. By the choice of
$\ell_*$, condition \eqref{cc:uf:trunc:eq:slow} holds, and
the lemma gives
$\|v\|_{[n],\gamma}\le C_n\mathsf D^3(\gamma(t_0)^{-1}\|v(t_0)\|_{[n]}
+\|\psi\|_{[n],\gamma})$.
The estimate divides by $\gamma$ after differentiation in label and
phase. Since $V=A_\perp v$, $v=A_\perp^\dagger V$ and
$\psi=A_\perp^\dagger f$, the product estimate
\eqref{cc:uf:trunc:eq:product} and
Lemma~\ref{cc:uf:trunc:lem:coefficients}(c) for $A_\perp$ and
$A_\perp^\dagger$ give \eqref{cc:uf:tower:eq:seed-norm}.

In (ii) the forcing vanishes for $t<t_0$, so $V(t_0)=0$.
Applying \eqref{cc:uf:tower:eq:seed-norm} to
\eqref{cc:uf:hier:eq:trans-ode} bounds $V$. By
\eqref{cc:uf:hier:eq:pressure-from-trans}, $\lambda\mathsf p$ is the
phase primitive of a product of the coefficients $\kappa/\lambda$ and
$\lambda^2/|\kappa|^2$ with $f-2MV$, so
\eqref{cc:uf:trunc:eq:phase-ops}, \eqref{cc:uf:trunc:eq:product} and
\eqref{cc:uf:hier:eq:coefficient-bound} give
\begin{equation}\label{cc:uf:hier:eq:pressure-bound}
 \lambda\|\mathsf p\|_{[n],\gamma}
 \le C_n\mathsf D^C\bigl(\|f\|_{[n],\gamma}+\|V\|_{[n],\gamma}\bigr),
\end{equation}
which proves \eqref{cc:uf:hier:eq:trans-bound}.

For (iii), the time independence of the envelope implies that $g\mathbf a$
solves the seed equation on $[t_0,T]$ at every label. It is perpendicular
to $\kappa$ and vanishes outside $\supp g$. By
\eqref{cc:uf:tower:eq:start}, its value at
the time $t_0$ is $g(\Id-P_\kappa)\mathbf a(t_0,0)$. The chain rule
gives $|g|_m\le C_m\ell_*^{-m}$, so $\|g\|_{[n]}\le C_n$. With
\eqref{cc:uf:hier:eq:coefficient-bound} for
$P_\kappa$ and \eqref{cc:uf:trunc:eq:product}, this gives
$\|g\mathbf a(t_0,\cdot)\|_{[n]}
\le C_n\mathsf D^C|\mathbf a(t_0,0)|=C_n\mathsf D^C\gamma(t_0)\lambda^{-3/4}$.
Hence \eqref{cc:uf:tower:eq:seed-norm} with $f=0$ gives
$\|g\mathbf a\|_{[n],\gamma}\le C_n\mathsf D^C\lambda^{-3/4}$.
For a function of the phase alone,
$\|F\|_{[n]}=\sum_{j\le n}h^j\|\partial_s^jF\|_\infty$, so
Lemma~\ref{cc:uf:exact:lem:profile} gives $\|F_h\|_{[n]}\le C_nh$.
By \eqref{cc:uf:trunc:eq:product},
$\|A_1\|_{[n],\gamma}\le C_n\mathsf D^C h\lambda^{-3/4}
=C_n\mathsf D^C\lambda^{-5/4}$. By \eqref{cc:uf:hier:eq:leading}, on
$[t_0,T]$ the leading oscillation solves \eqref{cc:uf:hier:eq:trans}
with $f=0$ and $\mathsf p=\mathsf p[A_1]$, so
\eqref{cc:uf:hier:eq:pressure-bound} gives \eqref{cc:uf:hier:eq:size}.
\end{proof}

\subsection{The finite expansion}\label{cc:uf:sec:hierarchy}
At depth $P$ we seek a perturbation
\[
 W=\sum_{p=1}^P Z_p+C_P,\qquad
 Z_p=Y_p+B_p,\qquad Y_p=V_p+C_{p-1},\qquad C_0=0,
\]
where $V_p$ is perpendicular to $\kappa$ and $B_p$ is independent of phase.
The divergence correction $C_p=C[V_p]$ makes
$V_p+C_p=\widetilde{\mathbf d}\times Q[V_p]$ divergence free by
\eqref{cc:uf:hier:eq:slaving}, which accounts for the last term $C_P$.
We call $p$ the \emph{level} of $Z_p$. The factor $|\kappa|^{-1}$
in the vector potential gives $C_{p-1}$ the bound of a term of
level $p$, so we group it with $V_p$.

\subsubsection{The equations at each level}

Start from $V_1,\mathsf p_1,C_1$ in
\eqref{cc:uf:hier:eq:switched} and put $B_1=0$, so that $Y_1=Z_1=V_1$.
Set aside the activation error $R_{\widetilde\chi}$ of
\eqref{cc:uf:hier:eq:switch}. The terms left by the leading wave at level
two are
\begin{equation}\label{cc:uf:hier:eq:H2}
 H_2=\widetilde\chi\bigl[(\partial_t+M)C[A_1]+\mathbf d\mathsf p[A_1]\bigr]
   +(V_1\cdot\mathbf d)V_1+(C_1\cdot\kappa)\partial_sV_1 .
\end{equation}
The linear term has zero phase average because $C[A_1]$ and
$\mathsf p[A_1]$ do. In the phase average of the last term,
\eqref{cc:uf:hier:eq:slaving} and integration by parts over a period give
\[
 \av{(C_1\cdot\kappa)\partial_sV_1}
 =-\av{\partial_s(C_1\cdot\kappa)V_1}
 =\av{(\mathbf d\cdot V_1)V_1}.
\]
Thus both quadratic terms in $\av{H_2}$ have a single label
derivative and two factors of $V_1$. Each factor has size
$C_n\mathsf D^C\lambda^{-5/4}$ by \eqref{cc:uf:hier:eq:size}, and the
label derivative costs a power of $\mathsf D$ by \eqref{cc:uf:trunc:eq:deriv}.
Hence $\|\av{H_2}\|_{[n]}\le C_n\mathsf D^C\lambda^{-5/2}$.
The parent residual $R_U\circ X$ satisfies a bound of this same order by
\eqref{cc:uf:hier:eq:parent-residual}, so it can be included in the
first mean equation.

In general, a product $(Z_i\cdot\mathbf d)Z_j$ has level $i+j$.
The fast derivative costs one level, so $(Z_i\cdot\kappa)\partial_sY_j$
has level $i+j-1$. Since $V_i\cdot\kappa=0$, only $C_{i-1}$ and $B_i$ can
differentiate the phase. We collect the terms involving
$(B_i\cdot\kappa)\partial_sY_j$ after the mean solve. The other
terms at level $p\ge3$ give
\begin{align}
 H_p={}&(\partial_t+M)C_{p-1}+\mathbf d\mathsf p_{p-1}\nonumber\\
 &+\sum_{i+j=p}(Z_i\cdot\mathbf d)Z_j
   +\sum_{\substack{i+j=p+1\\i\ge2}}
                  (C_{i-1}\cdot\kappa)\partial_sY_j,\label{cc:uf:hier:eq:Hp}
\end{align}
where all summation indices are positive integers.
For every $p\ge2$ put
\begin{equation}\label{cc:uf:hier:eq:source}
 f_p=-(H_p-\av{H_p}),\qquad \bar f_p=-\av{H_p}\quad(p\ge3),\qquad
 \bar f_2=-\av{H_2}-R_U\circ X.
\end{equation}
Every profile in $H_p$ is known from lower levels.
Apply \eqref{cc:uf:hier:eq:mean} to $\bar f_p$, obtaining $B_p,\bar q_p$.
Then put
\begin{equation}\label{cc:uf:hier:eq:supp}
 f_p^S=-\sum_{i=2}^p(B_i\cdot\kappa)\partial_sY_{p+1-i},
\end{equation}
and apply \eqref{cc:uf:hier:eq:trans} to $f_p+f_p^S$, obtaining $V_p,\mathsf p_p$.
Finally set $Q_p=Q[V_p]$, $C_p=C[V_p]$.
The term $-(B_p\cdot\kappa)\partial_sV_1$ in
\eqref{cc:uf:hier:eq:supp} accounts for the transport of the leading wave
by the new mean, so $B_p$ must be found before $V_p$. Every term in
\eqref{cc:uf:hier:eq:supp} has zero phase average since $B_i$ and
$\kappa$ are phase independent.

Since $D_aX$, $M$ and $\kappa$ are even in $a$, $R_U\circ X$ is odd in
$a$ and $A_1=g\mathbf aF_h$ changes sign under the reflection
$(a,s)\mapsto(-a,-s)$, the construction is symmetric under this
reflection: every $V_p$, $C_p$ and $B_p$ changes sign under it, and
$W^\natural$ is odd because $X$ is odd.

\begin{proposition}\label{cc:uf:hier:prop:identity}
For every integer \(P\ge2\), the profiles obtained from
\eqref{cc:uf:hier:eq:mean}, \eqref{cc:uf:hier:eq:source} and
\eqref{cc:uf:hier:eq:supp} are well defined, and
\[
 W=\sum_{p=1}^PZ_p+C_P=\sum_{p=1}^P(V_p+C_p+B_p),\qquad
 \mathsf P=\sum_{p=1}^P(\mathsf p_p+\bar q_p),\quad \bar q_1=0,
\]
satisfy \(\widetilde {\mathbf d}\cdot W=0\) and
\begin{equation}\label{cc:uf:hier:eq:residual}
 W_t+MW+(W\cdot\widetilde {\mathbf d})W+\widetilde {\mathbf d}\mathsf P
       +R_U\circ X=R_P+R_{\widetilde\chi}.
\end{equation}
Since $\av{V_p}=\av{C_p}=0$, the oscillatory part
\(W_{\rm osc}=W-\av W=\sum_{p=1}^P(V_p+C_p)\) is
supported in \(\supp g\), and so are the sources $f_p+f_p^S$.
The fields $V_p$, $C_p$, $\mathsf p_p$ and $f_p+f_p^S$ vanish for
$t<t_0$.
Set \(V_{P+1}=B_{P+1}=0\), so that \(Y_{P+1}=Z_{P+1}=C_P\). The
remainder of the expansion is
\begin{align}\label{cc:uf:hier:eq:tail}
 R_P={}&(\partial_t+M)C_P+\mathbf d\mathsf p_P\\
 &+\sum_{\substack{i,j\le P+1\\i+j\ge P+1}}
                  (Z_i\cdot \mathbf d)Z_j
 +\sum_{\substack{i,j\le P+1\\i+j\ge P+2}}
                  (Z_i\cdot\kappa)\partial_sY_j.\nonumber
\end{align}
At the labels outside $\supp g$ it is independent of the phase.
The construction is jointly smooth.
\end{proposition}
\begin{proof}
By \eqref{cc:uf:hier:eq:slaving} and \eqref{cc:uf:hier:eq:mean},
\[
 \widetilde{\mathbf d}\cdot W
 =\sum_{p=1}^P\bigl[\widetilde{\mathbf d}\cdot
   (\widetilde{\mathbf d}\times Q_p)+\mathbf d\cdot B_p\bigr]=0.
\]
With $W=\sum_{i=1}^{P+1}Z_i$, the quadratic term is
\[
 (W\cdot\widetilde{\mathbf d})W
 =\sum_{i,j}(Z_i\cdot\mathbf d)Z_j
  +\sum_{i,j}((C_{i-1}+B_i)\cdot\kappa)\partial_sY_j
\]
with $1\le i,j\le P+1$. For $2\le p\le P$, the terms to be
canceled at level $p$ sum to $H_p-f_p^S$, together with $R_U\circ X$
at $p=2$. The new linear terms
\[
 (\partial_t+M)(V_p+B_p)+\kappa\partial_s\mathsf p_p
    +\mathbf d\bar q_p=f_p+f_p^S+\bar f_p
\]
cancel them, since the right side is $-H_p+f_p^S$ for $p\ge3$ and
$-H_2+f_2^S-R_U\circ X$ for $p=2$. For $V_1$, \eqref{cc:uf:hier:eq:leading}
and \eqref{cc:uf:hier:eq:switched} give
\[
 V_{1,t}+MV_1+\kappa\partial_s\mathsf p_1=\widetilde\chi'A_1,\qquad
 (\partial_t+M)C_1+\mathbf d\mathsf p_1
 =\widetilde\chi\bigl[(\partial_t+M)C[A_1]+\mathbf d\mathsf p[A_1]\bigr]
  +\widetilde\chi'C[A_1].
\]
The two terms with $\widetilde\chi'$ sum to $R_{\widetilde\chi}$, and the
remaining term is canceled in $H_2$. The surviving
products have $i+j>P$ for $\mathbf d$ and $i+j-1>P$ for
$\kappa\partial_s$. Together with the two linear terms left at
level $P+1$, these give \eqref{cc:uf:hier:eq:tail}.

The sources at each level are smooth when the preceding profiles are.
The transverse equation is a linear ordinary differential equation at
fixed $(a,s)$ with smooth coefficients, the mean equation has smooth
solutions by Theorem~\ref{cc:uf:mean:inverse}, and taking the curl
preserves smoothness. By induction on the level, the profiles are well
defined and jointly smooth.

For the supports, let $\mathcal O$ be the open set of pairs $(t,a)$ with
$t<t_0$ or $a\notin\supp g$. The fields
$V_1=\widetilde\chi g\mathbf a F_h$, $C_1$ and $\mathsf p_1$ vanish on
$\mathcal O$, and so does the first term of $H_2$. If the oscillatory
profiles and pressures of index less than $p$ vanish on $\mathcal O$, then
on $\mathcal O$ the field $H_p$ contains only the phase-independent
products $(B_i\cdot\mathbf d)B_j$, and $f_p=f_p^S=0$ there. The solution
of \eqref{cc:uf:hier:eq:trans} then vanishes on $\mathcal O$, by the two
properties recorded after \eqref{cc:uf:hier:eq:trans-ode}, and so do
$C_p=\mathbf d\times Q[V_p]$ and, by
\eqref{cc:uf:hier:eq:pressure-from-trans}, $\mathsf p_p$, the zero-mean
primitive of a function that vanishes on $\mathcal O$. This proves the
assertions on the supports. On $\mathcal O$ the remainder \eqref{cc:uf:hier:eq:tail}
therefore consists of the products $(B_i\cdot\mathbf d)B_j$, which are
independent of the phase. The means $B_p$ do not vanish before $t_0$
in general.
\end{proof}

\subsubsection{Bounds at every order}

To bound the linear terms in \eqref{cc:uf:hier:eq:H2} and
\eqref{cc:uf:hier:eq:Hp}, we use the transverse equation to evaluate
the time derivative of the curl correction.
\begin{lemma}\label{cc:uf:hier:lem:relay}
There is an absolute exponent $C$ such that, for every \(n\ge0\) and
every field $V$ with $\av V=0$ and $\kappa\cdot V=0$,
\begin{equation}\label{cc:uf:hier:eq:completion-bound}
 \|C[V]\|_{[n],\gamma}\le C_n\mathsf D^C\lambda^{-1}\|V\|_{[n+1],\gamma}.
\end{equation}
If moreover \(V_t+MV+\kappa\partial_s\mathsf p=f\) with $\av f=0$ and
$\av{\mathsf p}=0$, then
\begin{equation}\label{cc:uf:hier:eq:relay-bound}
 \|(\partial_t+M)C[V]+\mathbf d\mathsf p\|_{[n],\gamma}
 \le C_n\mathsf D^C\lambda^{-1}
             \bigl(\|V\|_{[n+1],\gamma}+\|f\|_{[n+1],\gamma}\bigr).
\end{equation}
For the leading oscillation $V=A_1$ on $[t_0,T]$, the source in this estimate is
$f=0$.
\end{lemma}
\begin{proof}
By \eqref{cc:uf:hier:eq:completion}, $Q[V]$ is $\lambda^{-1}$ times the
phase primitive of $V$, multiplied by the coefficients $\kappa/\lambda$
and $\lambda^2/|\kappa|^2$, and $C[V]=\mathbf d\times Q[V]$.
Since $\ell_*^{-1}$ is a power of $\mathsf D$,
\eqref{cc:uf:trunc:eq:product}, \eqref{cc:uf:trunc:eq:deriv},
\eqref{cc:uf:trunc:eq:phase-ops} and
\eqref{cc:uf:hier:eq:coefficient-bound} give
\eqref{cc:uf:hier:eq:completion-bound}.

For \eqref{cc:uf:hier:eq:relay-bound} use \(\kappa_t=-M^T\kappa\),
\((|\kappa|^{-2})_t=2(\kappa\cdot M\kappa)|\kappa|^{-4}\),
and \([\partial_t,\mathbf d_j]=-M_{\ell j}\mathbf d_\ell\).
Differentiating $Q[V]$ in time and inserting the equation for $V_t$, in
which the term \(\kappa\times\kappa\,\partial_s\mathsf p\) vanishes, gives
\begin{equation}\label{cc:uf:hier:eq:potential-time}
 \partial_tQ[V]=\partial_s^{-1}\left[
  \frac{(M^T\kappa)\times V+\kappa\times MV-\kappa\times f}
       {|\kappa|^2}
  -2\frac{(\kappa\times V)(\kappa\cdot M\kappa)}
             {|\kappa|^4}\right].
\end{equation}
Commuting $\partial_t$ past the curl gives
$(\partial_t+M)C[V]=\mathbf d\times\partial_tQ[V]
+[\partial_t,\mathbf d\times]Q[V]+MC[V]$.
By \eqref{cc:uf:hier:eq:potential-time} and the formula for
$\partial_s\mathsf p$ in \eqref{cc:uf:hier:eq:pressure-from-trans},
each term, including $\mathbf d\mathsf p$, carries a factor $\lambda^{-1}$
and at most one label derivative. The product and primitive estimates
therefore give \eqref{cc:uf:hier:eq:relay-bound} without a time derivative
of $M$, $f$ or a mean. The weight $\gamma$ is applied after the time
differentiation.
\end{proof}

Up to fixed powers of $\mathsf D$, the leading oscillation has size
$\lambda^{-5/4}$ by \eqref{cc:uf:hier:eq:size}, its divergence correction
has size $\lambda^{-9/4}$ by \eqref{cc:uf:hier:eq:completion-bound},
and the first mean has size $\lambda^{-5/2}$ by the calculation after
\eqref{cc:uf:hier:eq:H2} and \eqref{cc:uf:mean:bound}.
A label derivative costs a power of $\mathsf D$, while a phase derivative
in Eulerian advection costs $|\kappa|h^{-1}\le\mathsf D^2\lambda^{3/2}$.
The smallest gain in frequency comes from advection of the phase by the divergence
correction, which contributes $\lambda^{-9/4}\lambda^{3/2}=\lambda^{-3/4}$.
Advection of the phase by the mean gives $\lambda^{-1}$,
as does taking a divergence correction. Advection in the labels gives
$\lambda^{-5/4}$. Allowing $\lambda^{-2/3}$ per level leaves a factor
$\lambda^{-1/12}$ to absorb the parent bounds.

We induct on the level at all derivative orders at once. When a source
requires a higher derivative of a previously computed profile, that
order is already available. The new mean $B_p$ enters the transverse
source \eqref{cc:uf:hier:eq:supp} undifferentiated.

\Needspace{10\baselineskip}
\begin{theorem}\label{cc:uf:hier:thm:depth}
There is an absolute exponent $C$ such that the following holds if
$\mathsf D^C\lambda^{-1/12}\le1$. For all integers \(n\ge0\) and
\(p\ge2\), the construction satisfies
\begin{align}
 \|V_p\|_{[n],\gamma}+\|f_p+f_p^S\|_{[n],\gamma}
        +\lambda\|\mathsf p_p\|_{[n],\gamma}
 &\le C_{n,p}\mathsf D^C\lambda^{-5/4}\lambda^{-2(p-1)/3},
 \label{cc:uf:hier:eq:Vbound}\\
 \|B_p\|_{[n]}+\|\bar f_p\|_{[n]}
 &\le C_{n,p}\mathsf D^C\lambda^{-5/2}\lambda^{-2(p-2)/3},
 \label{cc:uf:hier:eq:Bbound}\\
 \|C_p\|_{[n],\gamma}
 &\le C_{n,p}\mathsf D^C\lambda^{-9/4}\lambda^{-2(p-1)/3}.
 \label{cc:uf:hier:eq:Cbound}
\end{align}
The bounds \eqref{cc:uf:hier:eq:Vbound} and
\eqref{cc:uf:hier:eq:Cbound} also hold for $p=1$, with $f_1=f_1^S=0$.
For all integers \(P\ge2\) and \(m\ge0\), the expansion of
Proposition~\ref{cc:uf:hier:prop:identity} at the depth $P$ satisfies
\begin{equation}\label{cc:uf:hier:eq:sum}
 \|W_{\rm osc}\|_{[m],\gamma}\le C_{m,P}\mathsf D^C\lambda^{-5/4},\qquad
 \Big\|\sum_{p=2}^PB_p\Big\|_{[m]}+\Big\|\sum_{p=2}^P\bar f_p\Big\|_{[m]}
 \le C_{m,P}\mathsf D^C\lambda^{-5/2},
\end{equation}
and its remainder satisfies
\begin{equation}\label{cc:uf:hier:eq:tail-est}
 \|R_P\|_{[m]}\le C_{m,P}\lambda^{-5/4}\lambda^{-2P/3}.
\end{equation}
\end{theorem}
\begin{proof}
By \eqref{cc:uf:trunc:eq:product}, \eqref{cc:uf:trunc:eq:deriv} and
\eqref{cc:uf:hier:eq:coefficient-bound}, since $\ell_*^{-1}$ is a power
of $\mathsf D$ and $\kappa\partial_s$ costs
$|\kappa|h^{-1}\le\mathsf D^2\lambda^{3/2}$ by
\eqref{cc:uf:tower:eq:wave-vector}, fields $Z$, $Z'$, $Y$
satisfy, with the weight on either factor of a product or on neither,
\begin{equation}\label{cc:uf:hier:eq:advection}
 \begin{aligned}
 \|(Z\cdot\mathbf d)Z'\|_{[n],\gamma}+\|(\mathbf d\cdot Z)Z'\|_{[n],\gamma}
 &\le C_n\mathsf D^C\|Z\|_{[n+1]}\|Z'\|_{[n+1],\gamma},\\
 \|(Z\cdot\kappa)\partial_sY\|_{[n],\gamma}
 &\le C_n\mathsf D^C\lambda^{3/2}\|Z\|_{[n]}\|Y\|_{[n+1],\gamma}.
 \end{aligned}
\end{equation}
For $p=1$, \eqref{cc:uf:hier:eq:size} and
\eqref{cc:uf:hier:eq:completion-bound} give
\eqref{cc:uf:hier:eq:Vbound} and \eqref{cc:uf:hier:eq:Cbound} for
$A_1$, $\mathsf p[A_1]$, $C[A_1]$ on $[t_0,T]$, and
$0\le\widetilde\chi\le1$ gives them for $V_1$, $\mathsf p_1$, $C_1$.

To keep powers of $\mathsf D$ from accumulating with the depth, we
propagate a bound for the transverse source and the mean that contains
only powers of $\lambda$. For $p\ge2$ the induction claim at every order is
\begin{equation}\label{cc:uf:hier:eq:induction}
 \|f_p+f_p^S\|_{[n],\gamma}+\lambda^{1/2}\|B_p\|_{[n]}
 \le C_{n,p}\lambda^{-5/4}\lambda^{-2(p-1)/3}.
\end{equation}
For $B_p$ this is \eqref{cc:uf:hier:eq:Bbound} with $\lambda^{1/12}$ in
place of $\mathsf D^C$. Since \eqref{cc:uf:hier:eq:induction} contains
no power of $\mathsf D$, every exponent of $\mathsf D$ below is a sum
of a bounded number of the exponents of \eqref{cc:uf:mean:bound},
\eqref{cc:uf:hier:eq:trans-bound}, \eqref{cc:uf:hier:eq:size},
\eqref{cc:uf:hier:eq:completion-bound},
\eqref{cc:uf:hier:eq:relay-bound} and \eqref{cc:uf:hier:eq:advection},
and it depends neither on the level nor on the depth. The theorem holds
with the largest of these exponents, since $\mathsf D\ge2$. In the rest
of the proof every bound holds up to a factor $C_{n,p}$ or $C_{m,P}$,
which we do not write.

The bound \eqref{cc:uf:hier:eq:induction} at a level $i\ge2$ gives
\eqref{cc:uf:hier:eq:Vbound} and \eqref{cc:uf:hier:eq:Cbound} for $V_i$,
$\mathsf p_i$ and $C_i$. Indeed, by Proposition~\ref{cc:uf:hier:prop:identity} the source
$f_i+f_i^S$ vanishes for $t<t_0$ and outside $\supp g$, so
\eqref{cc:uf:hier:eq:trans-bound} bounds $V_i$ and $\lambda\mathsf p_i$
by $\mathsf D^C$ times the right side of
\eqref{cc:uf:hier:eq:induction}, and
\eqref{cc:uf:hier:eq:completion-bound} then bounds $C_i$. For the bounds
at each level it remains to prove \eqref{cc:uf:hier:eq:induction} and
\eqref{cc:uf:hier:eq:Bbound}.

Let $p\ge2$, and suppose that \eqref{cc:uf:hier:eq:induction} holds at
the levels $2\le i<p$ and at every order. Then
\eqref{cc:uf:hier:eq:Vbound} and \eqref{cc:uf:hier:eq:Cbound} hold for
$V_i$ and $C_i$, $1\le i<p$. Recall $Y_i=V_i+C_{i-1}$ and $Z_i=Y_i+B_i$.
The right side of \eqref{cc:uf:hier:eq:Cbound} for $C_{i-1}$ and the
bound \eqref{cc:uf:hier:eq:induction} for $B_i$ are at most the right
side of \eqref{cc:uf:hier:eq:Vbound} for $V_i$. Since
$\gamma\le1$, it follows that, for $i<p$,
\begin{equation}\label{cc:uf:hier:eq:YZ}
 \|Y_i\|_{[n],\gamma}+\|Z_i\|_{[n]}
 \le\mathsf D^C\lambda^{-5/4}\lambda^{-2(i-1)/3}.
\end{equation}

The source $f_p$ in \eqref{cc:uf:hier:eq:source} uses only profiles of
index less than $p$. Subtracting the phase mean of $(Z_i\cdot\mathbf d)Z_j$
removes $(B_i\cdot\mathbf d)B_j$, and every other product has a
factor $Y_i$ or $Y_j$, to which we apply the $\gamma$ norm.
By \eqref{cc:uf:hier:eq:relay-bound} and \eqref{cc:uf:hier:eq:Vbound}
for $V_{p-1}$, the linear term of $H_p$ is at most
$\mathsf D^C\lambda^{-9/4}\lambda^{-2(p-2)/3}$; for $p=2$
this is the bound for $A_1$, with $f=0$ and $\widetilde\chi\le1$. By
\eqref{cc:uf:hier:eq:advection} and \eqref{cc:uf:hier:eq:YZ}, a product
$(Z_i\cdot\mathbf d)Z_j$ with $i+j=p$ is at most
$\mathsf D^C\lambda^{-5/2}\lambda^{-2(p-2)/3}$. By
\eqref{cc:uf:hier:eq:advection}, \eqref{cc:uf:hier:eq:YZ} and
\eqref{cc:uf:hier:eq:Cbound} for $C_{i-1}$, a product
$(C_{i-1}\cdot\kappa)\partial_sY_j$ with $i+j=p+1$ is at most
$\mathsf D^C\lambda^{3/2}\lambda^{-9/4}\lambda^{-5/4}\lambda^{-2(p-2)/3}
=\mathsf D^C\lambda^{-2}\lambda^{-2(p-2)/3}$. The right side of
\eqref{cc:uf:hier:eq:induction} is
$\lambda^{-23/12}\lambda^{-2(p-2)/3}$, so the three bounds are this
number times
\[
 \mathsf D^C\lambda^{-1/3},\qquad
 \mathsf D^C\lambda^{-7/12},\qquad
 \mathsf D^C\lambda^{-1/12},
\]
respectively. Thus $\|f_p\|_{[n],\gamma}$ is at most
$\mathsf D^C\lambda^{-1/12}$ times the right side of
\eqref{cc:uf:hier:eq:induction}.

In the mean source $\bar f_p$, integration by parts in $s$ removes the
fast derivative just as at level two. By \eqref{cc:uf:hier:eq:slaving},
for $i\ge2$ we have
\begin{equation}\label{cc:uf:hier:eq:fast-mean}
 \av{(C_{i-1}\cdot\kappa)\partial_sY_j}
                 =\av{(\mathbf d\cdot V_{i-1})Y_j}.
\end{equation}
The linear terms have zero phase average because $C_{p-1}$ and
$\mathsf p_{p-1}$ do. Products of one mean and one oscillation also
average to zero, since $B_i$ and $\kappa$ are phase independent and
$\av{Y_j}=0$. Thus the mean source is
\[
 -\bar f_p=
 \sum_{i+j=p}\av{(Y_i\cdot\mathbf d)Y_j}
 +\sum_{i+j=p}\av{(\mathbf d\cdot V_i)Y_j}
 +\sum_{i+j=p}(B_i\cdot\mathbf d)B_j,
\]
with the additional term $R_U\circ X$ at $p=2$. By
\eqref{cc:uf:hier:eq:advection} and \eqref{cc:uf:hier:eq:YZ}, the first
two sums are at most $\mathsf D^C\lambda^{-5/2}\lambda^{-2(p-2)/3}$.
By \eqref{cc:uf:hier:eq:advection} and
\eqref{cc:uf:hier:eq:induction} at the lower levels, the third is at
most $\mathsf D^C\lambda^{-7/2}\lambda^{-2(p-2)/3}$, which is smaller.
The term $R_U\circ X$ is at most $C_n\lambda^{-5/2}$, by
\eqref{cc:uf:hier:eq:parent-residual}. Hence
$\|\bar f_p\|_{[n]}\le\mathsf D^C\lambda^{-5/2}\lambda^{-2(p-2)/3}$
at every order. With this estimate at the order $n+1$,
\eqref{cc:uf:mean:bound} bounds $\|B_p\|_{[n]}$ by $\mathsf D^C$ times
as much. This proves \eqref{cc:uf:hier:eq:Bbound} at the level $p$, and
with the hypothesis $\mathsf D^C\le\lambda^{1/12}$ it gives
\eqref{cc:uf:hier:eq:induction} for $B_p$.

For $f_p^S$ in \eqref{cc:uf:hier:eq:supp}, the means
$B_i$, $2\le i\le p$, are now bounded by
\eqref{cc:uf:hier:eq:induction}, and $Y_{p+1-i}$ is bounded by
\eqref{cc:uf:hier:eq:YZ}. By \eqref{cc:uf:hier:eq:advection}, each
product $(B_i\cdot\kappa)\partial_sY_{p+1-i}$ is therefore at most
$\mathsf D^C\lambda^{-3/2}\lambda^{-2(p-1)/3}$, which is
$\mathsf D^C\lambda^{-1/4}$ times the right side of
\eqref{cc:uf:hier:eq:induction}. With the bound for $f_p$,
$\|f_p+f_p^S\|_{[n],\gamma}$ is at most $\mathsf D^C\lambda^{-1/12}$
times that right side, and the hypothesis gives
\eqref{cc:uf:hier:eq:induction} at the level $p$. This completes the
induction. The number of products in a level and the constants of the
lower levels at higher orders enter $C_{n,p}$.

Since $\lambda\ge1$, summing
\eqref{cc:uf:hier:eq:Vbound}--\eqref{cc:uf:hier:eq:Cbound} over
$p\le P$ gives \eqref{cc:uf:hier:eq:sum}. For the remainder,
recall $Y_{P+1}=Z_{P+1}=C_P$. By \eqref{cc:uf:hier:eq:Cbound} for
$C_P$, the bound \eqref{cc:uf:hier:eq:YZ} holds for $i\le P+1$. The
products in \eqref{cc:uf:hier:eq:tail} are therefore bounded as those
of $H_p$ and $f_p^S$ were: the products $(Z_i\cdot\mathbf d)Z_j$ with
$i+j=q$, and $(C_{i-1}\cdot\kappa)\partial_sY_j$ and
$(B_i\cdot\kappa)\partial_sY_j$ with $i+j=q+1$, are at most
$\mathsf D^C\lambda^{-1/12}$ times $\lambda^{-5/4}\lambda^{-2(q-1)/3}$,
and in \eqref{cc:uf:hier:eq:tail} they have $q\ge P+1$. By
\eqref{cc:uf:hier:eq:relay-bound} and \eqref{cc:uf:hier:eq:Vbound} for
$V_P$, the linear term
\((\partial_t+M)C_P+\mathbf d\mathsf p_P\) is at most
$\mathsf D^C\lambda^{-1/3}$ times $\lambda^{-5/4}\lambda^{-2P/3}$. The
products of two means are bounded in the
unweighted norm, and every other term in the weighted norm, which
dominates the unweighted one because $\gamma\le1$. With the hypothesis,
this is \eqref{cc:uf:hier:eq:tail-est}.
\end{proof}
